\begin{lemma}
\label{editorial-source-lemma-characterize-integral-element}
Let $\varphi : R \to S$ be a ring map. Let $y \in S$. If there exists a
finite $R$-submodule $M$ of $S$ such that $1 \in M$ and $yM \subset M$,
then $y$ is integral over $R$.
Conversely, if $y$ is integral, the original subalgebra
$\varphi(R)[y]\subset S$ is a finite $R$-submodule containing $1$
and stable under multiplication by $y$.
\end{lemma}

\begin{proof}
Keep $\varphi:R\to S$ for the original coefficient map and write
$m_y:M\to M$, $x\mapsto yx$, for the $R$-linear endomorphism.
By Lemma \ref{editorial-source-lemma-charpoly-module} there is a monic
$P(T)=T^d+\sum_{i=1}^d a_iT^{d-i}\in R[T]$ with $P(m_y)=0$.
Evaluation at the original element $1\in M$ gives in $S$
$$
P^\varphi(y)=y^d+\sum_{i=1}^d\varphi(a_i)y^{d-i}
=P(m_y)(1)=0.
$$
Conversely, such a monic equation expresses $y^d$ as
$-\sum_{i=1}^d\varphi(a_i)y^{d-i}$. Multiplying this equation by
each needed power of $y$ and using induction expresses every power
of $y$ in the $R$-span of $1,y,\ldots,y^{d-1}$.
That span is exactly $\varphi(R)[y]$, contains $1$ and is stable
under $m_y$. If $S$ is the zero ring the same conclusion holds with
the zero submodule; a degree-one equation may be used.
\end{proof}